\Needspace{24\baselineskip}\fancyhead[L]{\small KOMET VON 1807.}
\centerline{KOMET VON 1807.}

\begin{center}
\rule{0.5\linewidth}{0.4pt}\\[2pt]
Monatliche Correspondenz. Band XVI. S. 562–567. Dezember 1807.\\[2pt]
\rule{0.5\linewidth}{0.4pt}
\end{center}
\par\indent ---\ ---\ --- Dr. \textsc{Gauss}, der die Bahn keines am Himmel erscheinenden Fremdlings unbeachtet lässt, schrieb uns am 16. November. „Aus den Bremer und Lilienthaler Beobachtungen bis zum 5. November habe ich folgende parabolische Elemente des Kometen abgeleitet

\smallskip

\begin{center}\begin{tabular}{lr}
Durchgang duchs Perihel, 18. Sept. .\ .\ .\ .\ . & $19^{\text{u}}\;6^{\text{m}}$ in Paris \\
Länge des Perihels .\ .\ .\ .\ .\ .\ .\ .\ .\ .\ . & $271^\circ\;0'\;13''3$ \\
$\AscendingNode$ .\ .\ .\ .\ .\ .\ .\ .\ .\ .\ .\ .\ .\ .\ .\ . & $266\;38\;31.2$ \\
\multicolumn{2}{l}{beide vom scheinbaren Äquinoktium aus gezählt} \\
Log. des kleinsten Abstandes .\ .\ .\ .\ .\ . & $9.8114927$ \\
{}[Neigung der Bahn .\ .\ .\ .\ .\ .\ .\ .\ .\ . & $63^\circ\;12'\;35''6$] \\
\multicolumn{2}{l}{Rechtläufige Bewegung}
\end{tabular}\end{center}


\par\indent Diese Elemente unterscheiden sich nur geringfügig von denen, die ich schon früher aus einigen Berliner Beobachtungen berechnet hatte. Ebenso weichen sie nur wenig von denen ab, die Herr \textsc{Bessel} aus einer kürzeren Reihe von Beobachtungen abgeleitet hat. Von Göttingen aus werde ich Ihnen mehr darüber schreiben.“

\begin{flushright}
[\textsc{von Zach}.]
\end{flushright}

